\documentclass[11pt,a4paper]{report}
\usepackage{ai_security_bluebook}

% Normalize “smart quotes” in source to TeX quotes (keeps PDF strings clean).
\usepackage{newunicodechar}
\newunicodechar{“}{``}
\newunicodechar{”}{''}
\newunicodechar{‘}{`}
\newunicodechar{’}{'}

% =============================================================================
% Narrative: one core with two wings (notes, not printed)
% Core: Unified AI Security Control Plane (the primary object)
% Wing A: Decoupled Safety (system principle)  Wing B: Evaluation Framework (evidence + governance engine)
% =============================================================================

% =============================
% Metadata (English edition)
% =============================
\bluebooktitle{Enterprise AI Security Technical Report}
\bluebooksubtitle{Unified AI Security Control Plane, Decoupled Safety, and the Evaluation Framework}
\bluebookauthor{Yunhao Wang et al.}
\bluebookorg{Digital Trust / Tianmu / AI Security Research}
\bluebookversion{v1.0-draft}
\bluebookdate{\today}
\bluebookclassification{Internal Draft}
\bluebookkeywords{AI Security; Unified Control Plane; Decoupled Safety; Runtime Enforcement; Evaluation Framework}
\bluebookcovernote{This volume treats \ControlPlane{} as the primary object; \DecoupledSafety{} as the system-principle wing; and \EvalFramework{} as the evidence and governance-engine wing. The three form a closed loop that supports deployable, auditable, and defensible enterprise AI security. The institutional positioning of organizational carriers and governance remit is given in Chapter~5 (\EvalFramework{} and \emph{guarantee production authority}). XeLaTeX is recommended.}

\begin{document}
\makebluebookfrontmatter

\begin{bluebookabstract}
This technical report argues for an enterprise-grade \textbf{guarantee production system} (an \textbf{obligation closure machine}): within a unified versioned boundary, it closes the loop among obligation semantics, runtime judgments, enforcement actions, evidence recording, and release statements, continuously producing \textbf{deployable, auditable, replayable, and contractible bounded guarantees}. The control plane names the smallest system boundary within which enterprise security obligations can be produced, maintained, and defended—not by detector count but by enforcing shared \texttt{policy\_version} and \texttt{correlation\_id} across judgment, enforcement, and evidence.

To realize this object, we adopt a “one core, two wings” structure. The \textbf{core} is the \textbf{Unified AI Security Control Plane}, which closes four primitives—\textbf{judgment, execution, evidence, and release}—within one boundary. \textbf{Wing A} is \textbf{Decoupled Safety}: rather than adding an external safety add-on, it rewrites problems that cannot be globally audited in an open semantic space into instance-level contracts that are event-decidable, policy-evaluable, failure-degradable, and evidence-recordable. \textbf{Wing B} is the \textbf{Evaluation Framework}: not merely scoring, but an evidence chain compiler that organizes benchmark constitutions, acceptance contracts, multi-track evaluation disciplines, and two evidence packs into a single governance substrate so that measurement, regression, and release gating enter one decision language. The shared objective is not “a bigger safety platform,” but the ability for enterprises to make lawful, verifiable, and revocable safety claims about their AI systems under explicit premises.

Our core judgment is that enterprise AI security hinges on whether the enterprise owns a control system that continuously closes judgment, execution, evidence, and release—not on ``whether the model is more aligned'' in isolation. As enterprise AI evolves from a single model endpoint into composite systems with gateways, retrieval, tools, workflows, memory, multi-turn state, and audit processes, risk expands into a system-level tension among confidentiality, integrity, availability, controllability, auditability, composability, and utility preservation. If security remains delegated to larger models, local guardrails, or offline scores, obligations fracture between policies, online behaviors, evidence chains, and release artifacts. The purpose of the control plane is to transform that fracture into versioned policy, enforceable actions, joinable evidence, and publishable bounded contracts.

This technical report makes \textbf{bounded, deployment-oriented claims}---not a global safety proof for open-ended natural language interaction. A claim is valid only when its scope, preconditions, degradation semantics, deployment regime, and evidence basis are explicit and can be compiled into an acceptance contract and a release evidence pack. Accordingly, evaluation is written as an \textbf{evidence engine for semantic bifurcation with governance convergence}: distinct jurisdictions, languages, and obligation spaces retain independent semantics on a shared execution backbone and converge at a unified governance decision board into a risk profile, release recommendation, and audit artifacts. The enterprise operating model compresses \textbf{theory, system, benchmark, and release} into one artifact chain so that every release becomes a defensible, replayable, rollbackable bounded contract.

\end{bluebookabstract}

\executivesummary{
The core problem of enterprise AI security is how to continuously produce, maintain, and defend signable safety guarantees. This technical report treats the primary object as a \textbf{guarantee production system}: under a unified versioned boundary, enterprises close obligation semantics, runtime judgments, enforcement actions, evidence recording, and release decisions into a single managed object, enabling safety claims to be produced, verified, contracted, replayed, and governed. The \ControlPlane{} is the system boundary form of that object—not its complete explanation.

We organize this object as “one core, two wings.” The core is the \textbf{Unified AI Security Control Plane}, responsible for closing \textbf{judgment, execution, evidence, and release}. The first wing, \textbf{Decoupled Safety}, rewrites safety from an implicit property buried in weights and prompts into instance-level contracts over policy, event, decision, evidence, and fail-safe objects. The second wing, the \textbf{Evaluation Framework}, organizes benchmark constitutions, acceptance contracts, multi-track disciplines, and two evidence packs into an evidence engine, so measurement, regression, and release defense share one governance grammar. AI security is thus not “more local defenses,” but an enterprise system that continuously produces bounded guarantees.

This lens is necessary because enterprise AI deployment is no longer single-model completion; it is a composite system of gateways, retrieval, tools, workflows, memory, multi-turn state, and audit processes. Risk expands beyond harmful outputs into systemic tensions across confidentiality, integrity, availability, controllability, auditability, composability, and utility preservation. If security is delegated to bigger models, more offline alignment, or loosely-coupled guardrails, organizations cannot maintain consistency among policy, online behavior, evidence chains, and release materials. Obligations fracture into isolated scores, isolated logs, and isolated blocks—never becoming a governable release object.

At the system-principle level, the necessary path is not “one more safety layer,” but retreating to a \textbf{computable envelope}. \DecoupledSafety{} compresses open task space into finite events, finite decision semantics, a finite fail-safe action space, and finite evidence fields, so hot-path decisions can be deterministically evaluated, written into evidence, and contractibly contracted under conflict. Runtime enforcement no longer depends on a model’s introspection under pressure; policy, the kernel, the fail-safe object, and the evidence graph carry executable obligations.

Methodologically, enterprise AI security is modeled as a \textbf{dual-timescale risk-adaptive control system}. The fast timescale is the online hot path: within single-request and single-session budgets, the system evaluates events, dispatches actions, converges fail-safe behavior, and appends evidence. The slow timescale is the offline cold path: benchmarks, risk profiles, multi-track evaluation, release gating, and regression continuously update the policy version, detector packs, capability packs, and release boundaries. Evaluation becomes a \textbf{decision-signal compiler}; evidence and release gating are not terminal archiving but the mechanism that makes judgment and execution auditable, releasable, and regressable guarantees.

At the evaluation and governance layer, we reject collapsing distinct risk semantics into one score and adopt \textbf{semantic bifurcation with governance convergence}. Global robustness, jurisdictional compliance, language-specific risks, and organizational obligations coexist on a shared execution backbone, each with its own suites, aggregation logic, and report semantics; they converge at a unified governance board to output release decisions—ship, ship with reduced privileges, canary, or rollback—rather than a hidden weighted score. Track A and Track B discipline is therefore a hard constraint: Track A carries deployable conclusions; Track B carries mechanism upper bounds and stress tests; the two must not be mixed. Any external or executive guarantee must be supported by a Mechanism Evidence Pack and a Release Evidence Pack.

We also specify a \textbf{claim discipline / guarantee grammar}. Any safety statement that appears in release defense, external material, or contractual text must declare its scope, invariants, degradation semantics, preconditions, evidence basis, and deployment regime. Statements that cannot be compiled into an acceptance contract and evidence pack are not valid enterprise guarantees. The output is a bounded deployment contract: under explicit premises, what guarantees can be claimed; when premises break, what minimal service form must be entered, and how external claims must contract accordingly.

Last, the unified R\&D framework is an \textbf{operating program}: a single artifact chain \textbf{theory $\rightarrow$ system $\rightarrow$ benchmark $\rightarrow$ release}. Details appear in the appendix on the unified R\&D operating program.
}

% =============================
% Chapter 1
% =============================
\chapterstarter{Unified AI Security Control Plane: Problem Statement and System Boundary}{
This chapter establishes the primary object of this technical report: enterprise AI security as an \textbf{obligation closure machine}, with \ControlPlane{} as its minimal engineering boundary defined by shared \texttt{policy\_version} and join keys so judgment, execution, evidence, and release mutually reconstruct. Using falsifiable fracture claims, responsibility topology, and a four-part closure model, we tighten the boundary, threat table, and CIA+ objectives into \ControlPlane{} as the smallest system form capable of producing, maintaining, and defending bounded guarantees.
}{
\begin{itemize}
  \item Failure claims and problem statement; obligation–control-point fractures and four-part closure.
  \item Boundary characterization and closure criteria, non-goals / out-of-scope architectures; boundary and threat table; CIA+ objectives.
  \item Claim summary (non-goals, bounded guarantees; detailed in Chapter~5 with the evidence engine and deployment contract).
\end{itemize}
}

\section{Strategic context and the enterprise shift}
\ControlPlane{}. As enterprise AI evolves from a single model endpoint into composite systems including gateways, vector/document indexes, workflows, tool calls, human-in-the-loop steps, and multi-turn state, risk expands from “harmful output” to systemic concerns: confidentiality, integrity, availability, controllability, and defensibility of governance. Prompt injection crosses trust boundaries; retrieval poisoning and hidden instructions enter context; tools map language into high-impact side effects; orchestration state accumulates across turns; and when logs and decisions cannot share a single correlation object, accountability and release chains fracture. In this environment, traditional appsec add-ons, content-safety modules, or “stronger model alignment” are insufficient: they may optimize local KPIs, but they cannot produce a unified, versioned policy, enforceable runtime actions, and replayable evidence under one version surface.

\section{System boundary and responsibility split}
\ControlPlane{}. We fix two responsibility boundaries: (i) \emph{security judgment responsibility}—policy interpretation, risk scoring, allow/deny/rewrite decisions, rationale codes, and structured evidence; (ii) \emph{execution platform responsibility}—networking, scaling, orchestration \textsc{sla}, compute placement, and load management. The semantic anchor of \ControlPlane{} belongs to the former, not the latter: it neither assumes platform orchestration responsibilities nor reduces security to “use a bigger model” or “add more bolt-on components.” Any external object that crosses the prompt boundary—documents, retrieved chunks, tool contracts, \textsc{api} responses—must be treated as part of the attack surface and cataloged and constrained under the same \texttt{policy\_version}, event semantics, and evidence schema.

\noindent\textbf{Coupled \textsc{sla} surfaces.} Guarantees over synchronous latency (the service invariant S1 in Chapter~4) assume both execution-platform budgets \emph{and} control-plane budgets for evaluation and evidence append within the same request envelope. If either side breaches budget, the externally contractible statement must degrade consistently with Chapter~5 acceptance clauses rather than silently blaming ``the other subsystem.''

\begin{figure}[htbp]
  \centering
  \IfFileExists{figures/Fig1.png}{\includegraphics[width=0.9\textwidth]{figures/Fig1.png}}{%
    \fbox{\parbox{0.9\textwidth}{\vspace{22mm}\centering\textbf{Figure 1 placeholder: figures/Fig1.png}\par\vspace{22mm}}}%
  }
  \caption{Responsibility-boundary topology of the Unified AI Security Control Plane}
  \label{fig:responsibility-topology}
\end{figure}

\begin{problemstatement}
The problem statement of this technical report is whether \textbf{judgment, execution, evidence, and release} can be closed under one version surface and one correlation key in dynamic context, retrieval, tools, side effects, memory, and multi-tenant constraints. If the four live in different frames, the system may be locally “safe,” yet it cannot form an enterprise-defensible guarantee: online effective policy diverges from release material; audits cannot reconstruct the policy and retrieval/tool events that governed a decision; regression and production do not share the same event and evidence objects. \textbf{Thus the conclusion compresses to:} under this technical report, \ControlPlane{} is \textbf{the minimal engineering boundary} that instantiates an obligation closure machine: it compiles judgment into execution, deposits execution into evidence, and compiles evidence into release.
\end{problemstatement}

\section{Why the control plane is necessary: obligations, fracture points, and the primary object}\label{sec:why-control-plane-necessary}
\ControlPlane{}.

\subsection{A necessity sketch for obligation closure}
Enterprise assurance is externally stateable only if there exists a single version surface on which obligation semantics, runtime judgments, enforcement acts, recorded evidence, and release statements can be mutually reconstructed for the same request---otherwise ``local safety'' cannot compose into an auditable system guarantee. Any architecture that splits those artifacts across incompatible identifiers, schemas, or tracks therefore cannot satisfy \emph{obligation closure}: there is no joint object on which to define a bounded contract. The control plane is the minimal engineering boundary proposed in this technical report precisely because it is where those artifacts are forced to share \texttt{policy\_version}, event types, and join keys---not because it aggregates more detectors.

\noindent\textbf{Three obligation fractures (why closure fails without a control plane).} First, \textbf{tenant and data obligations} cannot be stably expressed by a single weight: different tenants, jurisdictions, and data classifications require obligations to land in versioned policy bindings and auditable associations; otherwise model upgrades silently rewrite effective policy without a regression object. Second, \textbf{retrieval and tool obligations} attach to retrieval slices, index generations, tool capability matrices, and approval paths; without a unified event surface, releases cannot reconstruct “which external text and which high-impact call were allowed at that time.” Third, \textbf{release obligations} require release materials to cite the same policy pack, model routing, detector pack, and necessary telemetry samples as production; without a unified evidence schema, legal, data governance, and accountability cannot share one release object. Only a control-plane-shaped boundary converges obligation definition, online evaluation, evidence deposition, and release defense under one catalog and version surface.

\section{Boundary characterization and four-part closure criterion}\label{sec:boundary-characterization}
\ControlPlane{}. This technical report treats \ControlPlane{} as the boundary within which four primitives—\textbf{judgment, execution, evidence, release}—must share \texttt{policy\_version}, compatible event schemas, and join keys so that obligations remain mutually reconstructible under dynamics (Section~\ref{sec:why-control-plane-necessary}). Engineering accepts \ControlPlane{} where four-part closure is required; alternatives split identifiers or schemas and therefore cannot sustain bounded deployment contracts without reconciliation overhead.

\noindent\textbf{Testable closure criteria (engineering, not metamathematics).} (i)~Every governed synchronous decision cites \texttt{policy\_version} and leaves a joinable \texttt{decision\_record\_id} when the acceptance contract demands one. (ii)~No two planes disagree on identifiers for the same request without an explicit reconciliation or downgrade path. (iii)~Release materials reference the same version surface as online enforcement for the scope they claim.

\section{Non-goals and out-of-scope architectures}
\ControlPlane{}. \emph{Corollaries} of the boundary criterion (Section~\ref{sec:boundary-characterization}): pure model-centric stacks that optimize generation \textsc{sla} without a shared regression surface for tenant, tool, retrieval, and release obligations, and uncorrelated security features without shared event/evidence identifiers, cannot sustain bounded deployment contracts. An enforced control-plane-shaped boundary is what requires identical \texttt{policy\_version} online and offline and yields differential replay via unified evidence objects.

\section{Reading conventions (summary)}
\ControlPlane{} / \EvalFramework{}. This report commits to \emph{deployable, measurable, and revocable} contracts under explicit premises (Chapter~5). Benchmark governance defines measurement conditions; the control plane owns judgment and policy evaluation; the execution environment owns \textsc{sla} but not safety semantics; evidence/release gating and jurisdictional semantics (Chapter~6) produce auditable release objects.

\section{Engineering rule for guarantee production: version surfaces and join conditions}
\ControlPlane{}. \textbf{Engineering rule:} any bounded guarantee that can be stated externally must be reconstructible as a join over the evidence surface. Minimum keys include: \texttt{policy\_version}, \texttt{correlation\_id}, \texttt{decision\_record\_id}, \texttt{rationale\_code}, and (if Tier-R2 is triggered) \texttt{approval\_ref}. If any key is missing, release defense cannot prove “what was in effect, what was executed, what was recorded,” and the statement is not defensible.

\section{Assets, adversary model, and operating assumptions}
\ControlPlane{} / \EvalFramework{}. Protected assets include: confidentiality of user and enterprise data; integrity of prompts, retrieved content, tool parameters and side effects; service availability within \textsc{sla}; and integrity of audit records. Adversaries include: adversarial prompting; retrieval-channel poisoning or leakage; tool misuse and tool-chain attacks; manipulation of session/workflow state; and abuse of the operations plane (e.g., hidden instructions inside documents). Insider threats are primarily mitigated through separation of duties, evidence immutability, and least-privilege tooling—not solely model alignment.

Operating assumptions include: policy can be expressed in machine-checkable fragments; enforcement hooks exist at gateway, retrieval, inference, tools, and logging; some residual risks are explicitly accepted via an \emph{acceptance contract}; and nondeterminism must be bounded by statistical test plans and online monitoring rather than ignored.

\paragraph{Control-plane compute exhaustion.} Beyond mis-joins, the control plane faces \emph{resource starvation}: adversarial load, deep detector cascades, or evidence-bus backpressure can exhaust CPU, memory, or append quotas. Treat starvation as a first-class threat: per-tenant budgets, tier timeouts and skip rules (Chapter~4), admission control on evidence paths, and conservative D-actions under deadline pressure---not silent bypass of policy evaluation.

\begin{threatmodel}
\begin{threattable}
Control plane / evidence bus & Mis-joins, forged or missing correlation, log tampering, backpressure exhaustion & Poisoned or dropped evidence rows; \texttt{decision\_record\_id} minting failure; partition of append path & False assurance claims, false denials of service, audit/replay break & Append-only or hash-chained logs, admission control, replication, integrity checks (Chapter~3); hot-path degradation (Chapter~4) \\
Prompt interface & Policy bypass or exfiltration across the instruction boundary & Jailbreaks, delimiter attacks, indirect prompt injection & Policy failure, data leakage, privilege abuse & Input governance, normalization, instruction boundary \\
Retrieval path & Corpus manipulation / exposure of sensitive slices & Poisoning, conflicting sources, hidden instructions in docs & Wrong/harmful answers, leakage, compliance failure & Retrieval security, lineage, segmentation, citation binding \\
Tool invocation & Privilege escalation via side effects & Tool-chain attacks, parameter forgery, approval ambiguity & Data loss, financial/ops impact, lateral movement & Tool policy, approval, sandboxing, rate limits \\
Model execution & Unsafe generation under benign intent & Alignment cliff, factual hallucination & Harmful content, unreliable automation & Output governance, grounding constraints, calibrated abstention \\
Runtime orchestration & Control-plane inconsistency & Cross-turn state manipulation, stale policy binding & Silent policy drift, audit breaks & Versioned policy engine, atomic decision logging \\
\end{threattable}
\end{threatmodel}

\section{Security objectives and trade-offs}
\ControlPlane{} / \EvalFramework{}. The objectives below must be jointly satisfied at the system level; each includes the criterion “if not in the control plane, it must fail,” to pin down the necessity of the primary object.

\begin{securityobjectives}
\begin{itemize}
  \item \textbf{Confidentiality.} Sensitive data must not cross policy trust boundaries via prompts, retrieval, logs, or tools. \textbf{If not in the control plane:} independent redaction modules cannot prove “the same retrieval–tool chain” did not leak obligation-bearing slices.
  \item \textbf{Integrity.} Prompts, retrieval chunks, tool parameters, and side effects are checkable under the threat model. \textbf{If not in the control plane:} release and online evaluation use different field definitions, making integrity claims indefensible.
  \item \textbf{Availability.} Service meets \textsc{sla} with measurable degradation. \textbf{If not in the control plane:} security and platform report latency inconsistently; a single “p99 with security enabled” contract cannot be signed.
  \item \textbf{Controllability.} High-risk behaviors are constrained by explicit runtime policy. \textbf{If not in the control plane:} “model habits” cannot be versioned or gated by regression.
  \item \textbf{Auditability.} Policy-relevant decisions emit replayable structured evidence. \textbf{If not in the control plane:} audit and production logs cannot share correlation IDs; postmortems degrade into guesswork.
  \item \textbf{Composability.} Controls remain effective across heterogeneous models and stacks. \textbf{If not in the control plane:} every model swap implies a new implicit policy; evidence cannot accumulate comparably.
  \item \textbf{Utility preservation.} Minimize needless harm; unnecessary blocks are defects. \textbf{If not in the control plane:} local modules optimize block rates while global utility has no owner.
\end{itemize}
\end{securityobjectives}

\noindent\textbf{Trade-offs.} Confidentiality/integrity often conflict with utility and availability; audit depth conflicts with cost and latency. These tensions must be \emph{policy-ized} inside the control plane (versions, thresholds, exceptions, evidence packs) rather than resolved by ad hoc cross-team negotiation. \DecoupledSafety{} (Chapter 2) supplies the externalized policy-first mechanism; \EvalFramework{} (Chapter 5) supplies measurable and releasable evidence—both depend on the boundary fixed in this chapter.

\designprinciple{Decoupling (to support judgment)}{Separate policy from model weights so “which rule applied at that time” is citeable and signable.}
\designprinciple{Cross-layer enforcement (to support execution)}{Enforcement hooks across input, retrieval, inference, tools, output, and orchestration logs share one event and version semantics.}
\designprinciple{Evidence before trust (to connect evidence and release)}{Any internal/external guarantee must point to tests, telemetry, or signed artifacts, under the same version surface as the release object.}

\keytakeaway{The central contradiction of enterprise AI security is whether the four primitives close: judgment, execution, evidence, and release must be joinable and replayable under the same catalog and version surface of \ControlPlane{}. Otherwise, obligations inevitably fracture in composite systems—therefore the primary object is \ControlPlane{} as the minimal boundary that carries shared version and join semantics; model platforms and uncorrelated feature stacks do not provide that closure by themselves.}

% =============================
% Chapter 2
% =============================
\chapterstarter{Decoupled Safety: Conceptual Foundations and Safety Model}{
\DecoupledSafety{} is a \emph{necessary} system path in enterprise AI: under open semantics, retrieval, tools, multi-turn state, and deployment constraints, safety must be written as an \emph{instance-level system contract} that is evaluable, degradable, and evidentiary. It compresses unsignable “semantic safety wishes” into hot-path event evaluation, action dispatch, and fail-safe protocols, and cold-path version-monotone policy and evidence updates—enabling guarantee production.
}{
\begin{itemize}
  \item From open semantics to instance contracts: necessity and the computable envelope.
  \item Three foundational failure modes and their design implications: what must be externalized.
  \item Minimal object model and bounded claims: shared object language and claim scope for later chapters.
\end{itemize}
}

\section{The necessary path: from open semantics to a computable envelope}
\textbf{Formal envelope (finite bases).} Let finite sets \(\mathcal{E}\) (event types), \(\mathcal{A}\) (actions including fail-safe), \(\mathcal{R}\) (rationale codes), and \(\mathcal{V}\) (policy artifact versions). Hot-path semantics are maps \(\mathcal{E}\times\mathcal{V}\rightarrow \mathcal{A}\times\mathcal{R}\) subject to declared budgets; anything outside these carriers is out-of-band until promoted via cold-path version bumps.

The first challenge of enterprise AI security is not “raise the model’s safety tendency by a bit,” but that in an open semantic space a system cannot perform global, complete, low-cost safety adjudication on the hot path. Open-ended tasks, multi-turn state, retrieval injection, tool side effects, and jurisdictional obligations make “complete online semantic review” infeasible and un-signable as an enterprise contract. The deployable alternative is therefore not to keep chasing global semantic review, but to retreat to a \textbf{computable envelope}: the system does not need to fully understand all natural language meaning at runtime; it must project the open space into finite event types, finite actions, finite degradation protocols, and finite evidence fields so that hot-path decisions are deterministically evaluable, written into evidence, and fail-safe contractible under conflict.

From this perspective, \DecoupledSafety{} is not merely “taking safety out of the model.” It rewrites safety from an unsignable semantic aspiration into an instance-level contract that can be versioned as \texttt{policy\_version}, typed as \texttt{event} schema, decided under finite \texttt{decision} semantics, and recorded into an evidence graph. The hot path carries \textbf{budgeted evaluation, dispatch, and evidence emission}; the cold path carries \textbf{versioned artifact updates without online semantic discontinuities}. This is the system premise for the runtime protocols, evaluation evidence packs, and bounded deployment contracts in later chapters.

\paragraph{Projection map and residual false negatives.} Let contexts in the open semantic space map to carriers in \(\mathcal{E}\) via a projection \(\pi\) (normalization + typing + discrete buckets). \(\pi\) is necessarily lossy: attacks or obligations that do not surface as discriminable \texttt{event} features become \emph{false negatives} at the hot-path boundary. Such residual error is not eliminable by stacking detectors alone; it must be bounded by cold-path evidence: Track~B stress suites, ORA-driven detector and policy artifact updates, and acceptance-contract slices that explicitly own the residual risk (Chapter~5). \textbf{Quantitative coupling:} operational \emph{loss rate} attached to \(\pi\) for a given obligation class is measured by false-negative and miss-rate indicators in Track~B suites (labeled slices, scenario coverage), feeding the same acceptance-contract error budgets that Chapter~5 uses for deployable vs.\ upper-bound claims.

\paragraph{Engineering semantics vs.\ formal proof.} The finite sets \(\mathcal{E},\mathcal{A},\mathcal{R},\mathcal{V}\) and maps in this chapter are \textbf{contractible engineering semantics}: they define what implementations and acceptance contracts must align to for release and audit. They are \emph{not} a claim of a complete formal correctness proof of an open-ended AI system.

\subsection{Compilable, partially compilable, and non-compilable obligations}
Obligations differ in how far they can be reduced to event-typed, versioned checks. The table below is a \emph{governance partition} for what may enter Track~A vs.\ what must stay Track~B, statistical, or post-release supervision (see also ``Compilable vs.\ non-compilable'' in Chapter~5).

\begin{center}
\footnotesize
\begin{tabular}{@{}p{0.30\textwidth}p{0.32\textwidth}p{0.32\textwidth}@{}}
\toprule
\textbf{Obligation class} & \textbf{Compilability} & \textbf{Admissible claim shape} \\
\midrule
Deterministic structural (allowlists, schema tags, regex, tool-destination rules) & \textbf{Event-typed:} directly bound to \texttt{event} / \texttt{policy} & Track~A mechanism + S1-style bounds when evidence joins \\
Probabilistic / tail risk (adaptive attacks, rare harm) & \textbf{Partial:} needs explicit error budget, slices, Track~B stress & Statistical or upper-bound claims; migration to Track~A only with contract + replay \\
Open-ended virtues (``always helpful,'' global harmlessness) & \textbf{Non-compilable} as unconditional hot-path guarantees & Research agenda, non-goals, or explicit statistical claims with scope \\
\bottomrule
\end{tabular}
\end{center}

\section{Introduction: the gap between alignment and system security}
Supervised fine-tuning, preference alignment, and safety post-training reduce harmful completions in isolation but do not imply enterprise system security for composite systems with tenant obligations, retrieval lineage, tools, and release constraints (Chapter~1). \DecoupledSafety{} is the complement, not the replacement, for alignment: generation stays with the model; policy, fail-safe protocols, evidence schema, and release logic are externalized as stable, versioned objects.

\section{Three foundational failures and their design implications}
\textbf{Failure 1: semantic review failure and compute–verification asymmetry.} In open tasks with rich context, setting “review all harmful semantics” as a unified online goal immediately conflicts with executability, compute budgets, latency, and test semantics. The design implication is not “add stronger reviewers,” but to contract verification to instance objects: discrete \texttt{event}s (input blocks, retrieval units, tool calls, candidate outputs) must be evaluable under versioned \texttt{policy}; the \texttt{kernel} must produce \texttt{decision}s with rationale codes and joinable \texttt{evidence}.

\textbf{Failure 2: cognitive consistency failure and the need for value/ontology anchoring.} Multi-turn dialogue, socio-technical pressure, and tool-chain side effects break the assumption “the model will remain safe after deeper reflection.” The implication is that value boundaries, ontology constraints, and the acceptable action space must be written into external \texttt{policy}. When conflict signals hold, the system must invoke the controlled degradation defined by \texttt{FailSafeObject}—not continue open-ended generation. Recovery must be authorized, evidentiary, and citeable.

\textbf{Failure 3: retrieval-path failure and the need for representation/behavioral monitoring.} Surface tokens and local prompt semantics cannot cover retrieval poisoning, hidden instructions, cross-document assembly, and ranking manipulation. The implication is that retrieval must enforce lineage, segmentation, and citation binding; monitoring and evidence must preserve representation-level or behavior-trace signal interfaces so “what the system saw, why it judged anomalous, why it degraded or reduced privileges” is capturable in the evaluation and evidence system of Chapter 5.

\section{Minimal object model (interfaces to Chapters 3--5)}
To avoid \DecoupledSafety{} collapsing into slogans, we fix the minimal object set that implementations and evidence must carry across Chapters 3--5. \texttt{policy} represents versioned obligations, thresholds, and allowed action spaces; \texttt{kernel} consumes \texttt{event}s to produce \texttt{decision}s and emit \texttt{evidence}; \texttt{environment} describes topology, trust zones, observation budgets, and deployment layering; \texttt{event} is a typed security event such as an input block, retrieval hit, tool call, or candidate output; \texttt{decision} is a finite action (allow/rewrite/sandbox/escalate/block/degrade) plus a rationale code; \texttt{evidence} is a joinable, replayable, differential record set; \texttt{FailSafeObject} binds trigger conditions, allowed degradation sets, service invariants, and recovery tiers into one protocol object; if out-of-band adaptation is enabled, \texttt{OOB\_adaptation} describes the cold-path configuration, artifacts, and evidence anchors.

\section{Typical failure patterns}
\securityrisk{Attribution ambiguity}{When safety logic is hidden in prompt tricks or model weights, upgrades silently rewrite effective policy without a regression object.}
\securityrisk{Orchestration bypass}{Agents and workflows route through alternate tools or hidden channels to bypass enforcement, collapsing control consistency.}
\securityrisk{Evidence absence}{Without a stable evidence graph, the organization cannot reconstruct “why it allowed/blocked, under which policy version, with which retrieval/tool context.”}

\section{Bounded claims and the safety contract of this volume}
\DecoupledSafety{} is bounded and revocable: it does not claim global, complete, decidable safety for arbitrary language tasks and tool ecosystems, nor does it claim models never err. It claims that \textbf{under explicit environment disclosure, versioned policy and detector artifacts, measurable observation and compute budgets, and a signed acceptance contract}, the system can raise attacker cost, limit blast radius, and enter a pre-declared minimum-safe service form when premises break.

\researchagenda{Foundational question}{Under heterogeneous models, nondeterministic decoding, and multi-step tool use, which compositional constraints can be enforced within a measurable error budget while evidence packs still satisfy Track A (deployable) release discipline?}

\keytakeaway{\DecoupledSafety{} is not an add-on layer. It is the system principle that rewrites safety into a computable envelope, instance contracts, fail-safe protocols, and joinable evidence—enabling enterprise AI security to be produced, maintained, audited, and defended in release.}

% =============================
% Chapter 3
% =============================
\chapterstarter{Reference Architecture and the Decoupled Safety Kernel}{
This chapter provides a deployable \emph{system-level} view: layering remains the skeleton, but we must also describe \emph{control flow and evidence flow}, the \textbf{minimal fields of three planes}, the split between \textbf{synchronous online enforcement} and \textbf{asynchronous out-of-band adaptation}, and the \SafetyKernel{} minimal interface—so architecture diagrams are no longer static posters but the mount points for Chapter 4 protocols and Chapter 5 evidence.
}{
\begin{itemize}
  \item Layered responsibilities + dynamic control/evidence paths (and how they relate to static box diagrams).
  \item Data plane / control plane / evidence plane: minimal field sets and join rules.
  \item Synchronous enforcement vs asynchronous OOB adaptation; \SafetyKernel{} minimal interface table and deliverables.
\end{itemize}
}

\section{From responsibility topology to join topology: how boundaries are closed by evidence}
\ControlPlane{}. Chapter 1 defined the control plane as a responsibility boundary topology. At the system level, this must become a \textbf{join topology}: which keys must appear across the data/control/evidence planes to reconstruct a single narrative that binds one judgment to its execution, evidence, and release. Thus “reference architecture” is judged not by component count but by \textbf{closure capacity}: the same \texttt{policy\_version} constrains hot-path evaluation and cold-path aggregation; the same \texttt{correlation\_id} threads the event sequence; the same \texttt{decision\_record\_id} enables release gates to cite evidence. Missing any key collapses the control plane into decoration.

\section{Dual-plane correctness: hot-path atomic enforcement and cold-path version monotonicity}
\ControlPlane{}. The dual-plane split is not only performance—it is correctness. The \textbf{hot path} must complete deterministic evaluation, controlled dispatch, and \emph{atomic evidence appends} within S1 budgets; the \textbf{cold path} may update control laws only through versioned artifacts (\texttt{policy\_version}, detector packs, index generations, routing tables) and only after passing release gating. Any “online ad hoc change” that bypasses version monotonicity breaks replay and defensibility, violating the core promise of the guarantee production system.

\section{Layered reference architecture and dynamic control flow}
\ControlPlane{}. Layering answers “who owns what”; \emph{dynamic control flow} answers “how state transitions across one request.” Both are required: interaction–context–model–runtime–governance layers are connected at runtime by \textbf{events}. Under the same \texttt{correlation\_id}, the system emits context bindings, retrieval hits, model routing decisions, \SafetyKernel{} evaluations, and evidence appends. Static diagrams that do not align to the event sequence produce “the diagram exists but the logs do not” fractures in engineering.

\section{Data plane, control plane, and evidence plane: minimal fields}
\ControlPlane{}. The three planes are not philosophical labels; they are minimal column sets that must be written and validated. Implementations may add fields, but must not split join keys.

\begin{center}
\footnotesize
\begin{tabular}{@{}p{0.18\textwidth}p{0.78\textwidth}@{}}
\toprule
\textbf{Plane} & \textbf{Minimal fields (example names; implementations may map to schema)} \\
\midrule
Data plane & \texttt{correlation\_id}; \texttt{tenant\_id} / \texttt{principal\_id}; \texttt{session\_id}; \texttt{prompt\_digest} (or equivalent); \texttt{retrieval\_unit\_id[]} (incl. slice version / index generation); \texttt{tool\_invocation\_id[]} (incl. parameter digest); \texttt{model\_route\_id}; \texttt{completion\_id}; sensitivity labels (if applicable). \\
Control plane & \texttt{policy\_version}; \texttt{policy\_binding\_id}; \texttt{risk\_score} (or discrete risk tier); \texttt{intervention} (allow/rewrite/\ldots); \texttt{route\_decision} (e.g., light vs heavy detection path); \texttt{fail\_safe\_state} (if degraded); \texttt{deadline\_budget\_ms} (aligned with \textsc{sla}). \\
Evidence plane & Joinable \texttt{decision\_record\_id}; \texttt{evidence\_schema\_version}; \texttt{rationale\_code}; optional \texttt{artifact\_hash[]}; human approval pointer \texttt{approval\_ref} (if Tier-R2); immutable storage locator (e.g., object-store keys). \\
\bottomrule
\end{tabular}
\end{center}

\noindent\textbf{Join rule.} Any \textbf{decision} must be uniquely reconstructible in the evidence plane from a (data-plane subset + control-plane subset). Release gates and regression only recognize records satisfying this join closure (aligned with Chapter 5 minimum evidence packs).

\subsection{Join as a primary invariant: compositional evidence grammar}
Treat the join not as documentation glue but as a \textbf{primary system invariant}: over a single \texttt{correlation\_id}, component-level evidence must compose into one replayable narrative under one \texttt{policy\_version}. Write schematic component records $E_{\mathrm{in}}, E_{\mathrm{ret}}, E_{\mathrm{gen}}, E_{\mathrm{tool}}$ for gateway/input binding, retrieval lineage, generation/kernel evaluation, and tool side effects. A system-level claim uses a compositional join (same identifiers, mutually consistent \texttt{rationale\_code} semantics):
\[
  G_{\mathrm{sys}} \equiv E_{\mathrm{in}} \oplus E_{\mathrm{ret}} \oplus E_{\mathrm{gen}} \oplus E_{\mathrm{tool}} \quad\text{(all under shared \texttt{correlation\_id}, \texttt{policy\_version}).}
\]
\textbf{Legal compositions:} every node cites the same \texttt{policy\_version}; evidence rows join on declared keys; rationale codes do not contradict along the chain. \textbf{Illegal compositions (shadow bypass):} generation evidence marks the response ``safe'' while tool evidence is missing for a request that executed high-impact side effects; or retrieval lineage is absent while a sensitive slice informed the prompt---these invalidate system-level Track~A conclusions even if a component score looks good.

\begin{figure}[htbp]
  \centering
  \IfFileExists{figures/Fig3.1.png}{\includegraphics[width=0.92\textwidth]{figures/Fig3.1.png}}{%
    \fbox{\parbox{0.92\textwidth}{\vspace{18mm}\centering\textbf{Figure 3.1 placeholder: figures/Fig3.1.png}\par\vspace{4mm}\footnotesize Gateway $\rightarrow$ Retrieval $\rightarrow$ Kernel $\rightarrow$ Tool layers threaded by \texttt{correlation\_id}, producing joinable rows for the Release Evidence Pack.\par\vspace{18mm}}}%
  }
  \caption{Evidence-join dependency graph (conceptual): \texttt{correlation\_id} threads layers into a single release-relevant record set}
  \label{fig:evidence-join-graph}
\end{figure}

\noindent\textbf{Engineering closure principle (compositional evidence, informal).} A system-level deployable claim is \emph{defensible} only if every mandatory sub-component record for that \texttt{correlation\_id} is present, joinable, and semantically consistent within the declared synchronous budget S1; missing or mismatched joins refute the claim regardless of partial metrics.

\section{Joinable evidence bus: failures, partitions, and assurance--availability tension}
\ControlPlane{}. The compositional grammar $G_{\mathrm{sys}}$ assumes rows append reliably with shared keys. Distributed realities impose explicit modes:
\begin{itemize}
  \item \textbf{Evidence-path degradation.} If append latency exceeds S1, the synchronous API must still satisfy S1 using Chapter~4 fail-safe (typically conservative D-actions with explicit error envelopes). Assurance statements that require complete joins must \emph{contract}---they cannot silently pretend full evidence arrived.
  \item \textbf{Partition / outage of logging.} \emph{Fail-closed on assurance:} missing mandatory $E_{\mathrm{tool}}$ or $E_{\mathrm{ret}}$ when side effects occurred forbids Track~A ``safe'' verdicts until replay reconstructs rows or scope narrows. \emph{Fail-open on availability (explicit exception only):} an acceptance-contract clause may declare degraded observability regions where only weaker guarantees apply.
  \item \textbf{\texttt{decision\_record\_id} mint failure.} Without an identifier, hot-path execution must not proceed as if governed evidence exists; emit rationale codes tied to ``record withheld'' states and trigger Tier-R1/R2 per policy.
\end{itemize}

\paragraph{Hash-chained immutability vs.\ S1.} When \texttt{decision\_record\_id} issuance or append paths use hash-chained or append-only storage, the \textbf{durable chain extension} (replication, cross-partition ordering) is an asynchronous concern: the hot path must still return within S1 with a stable record handle and (if used) a hash preimage or log-sequence token that the client and release pack can cite. Full hash verification and backfill must not block the synchronous response; doing so would conflate proof latency with user-facing \textsc{sla} and should be budgeted on a separate durability path, with fail-safe behavior if only the fast append path is available (Chapter~4).

\paragraph{Canonical failure pattern (cross-tenant tool escalation via retrieval).} An adversary poisons or manipulates retrieval so that the model proposes a tool call that exfiltrates across tenant boundaries. If retrieval evidence $E_{\mathrm{ret}}$ does not bind slices to \texttt{decision\_record\_id}, or tool evidence $E_{\mathrm{tool}}$ is absent while calls occurred, the Release Evidence Pack must fail \emph{incomplete evidence} under admissibility rules (Chapter~5)---the gate rejects ``safe overall'' narratives built only on generation scores.

\noindent\textbf{Join completeness and operating mode.} Full join completeness is a \textbf{precondition for Track~A deployable closure}, not for every degraded-availability response: synchronous APIs may still return under weaker contract classes. \textbf{Mandatory-evidence bridge:} if mandatory rows are missing yet execution continues, it may do so only under acceptance-contract tiers that explicitly weaken assurance; Track~A ``safe'' narratives and strong release clauses are \textbf{suspended} until replay reconstructs joins or scope contracts (Chapter~5 admissibility).

\section{Synchronous online enforcement plane and asynchronous out-of-band adaptation plane}
\ControlPlane{} + \DecoupledSafety{}.

\paragraph{Synchronous online enforcement plane (hot path).} Within single-request/session \textsc{sla}, the system must complete event normalization, policy binding, deterministic evaluation, allowed action dispatch, synchronous response formation, and \emph{atomic evidence append}. Budgeted work may include light rules, cached classifiers, bounded retrieval checks, and one pre-warmed heavy detection pass. Principle: do \textbf{not} depend on cross-request global optimization, and do not block on human approval unless entering a protocolized Tier-R2 that still satisfies S1.

\paragraph{Asynchronous out-of-band adaptation plane (cold path).} Outside request boundaries: corpus/index refresh, drift aggregation, policy candidate mining, full regression replays, red-team campaigns, threshold/routing sweeps, and offline model-route evaluation. Outputs are \emph{next-version} artifacts—\texttt{policy\_version}, detector packs, index generations, routing tables—and only become effective in the hot path after release gating.

\paragraph{Separation discipline.} The hot path consumes only \emph{released} versions; the cold path must not mutate hot configurations without version monotonicity and approval. Fail-safe triggers are primarily hot-path signals; adaptations are produced on the cold path.

\section{\SafetyKernel{} minimal interface}
\DecoupledSafety{}. The following four calls form a minimal implementable set, corresponding to Chapter 2’s \texttt{event}/\texttt{decision}/\texttt{evidence}.

\begin{center}
\footnotesize
\begin{tabular}{@{}p{0.22\textwidth}p{0.32\textwidth}p{0.40\textwidth}@{}}
\toprule
\textbf{Interface} & \textbf{Inputs (summary)} & \textbf{Outputs (summary)} \\
\midrule
\texttt{ingest\_event} & raw or semi-normalized payload + tenant/session context & \texttt{event\_id} + normalized \texttt{event} record \\
\texttt{evaluate} & \texttt{event\_id} + \texttt{policy\_version} + optional feature-cache handle & \texttt{decision} + \texttt{rationale\_code} + risk tier \\
\texttt{dispatch} & \texttt{decision} + execution context & controlled downstream side effects / block / rewrite result \\
\texttt{emit\_evidence} & \texttt{decision} + data-plane snapshot pointer & \texttt{decision\_record\_id}, joinable evidence-plane record \\
\bottomrule
\end{tabular}
\end{center}

\noindent\textbf{Non-goals.} The kernel does not implement full business orchestration semantics; does not replace platform scheduling; and does not perform cold-path retraining—only guarantees hot-path enforcement and evidence atomicity.

\section{Controls-core: engineering instantiation of \SafetyKernel{}}
\DecoupledSafety{}. The controls-core submodule shows that \SafetyKernel{} can appear first as a minimal engineering instance without expanding immediately into a full control plane. Its abstractions consolidate into four object families—\textbf{types \& interfaces, detection, scoring, prevention}—with three build tiers (\textbf{controls-core-lite / controls-core-advanced / controls-core-full}) and two deployment carriers (\textbf{controls-local-service / controls-cloud-service}), forming the basic layering of capability packs and deployment regimes. In other words, an engineered kernel is defined not by “how many detectors exist,” but by whether there are stable event interfaces, tiered execution, switchable capability packs, and governable service carriers.

In this volume’s object language, this maps to minimal kernel engineering: types \& interfaces correspond to event and detector contracts; detection / scoring / prevention correspond to decision synthesis and controlled actions; lite / advanced / full correspond to capability tiers; local / cloud services correspond to concrete deployment carriers. The submodule also supports \textbf{registry auto-discovery, config enable/disable, API / worker pipelines, DB / object storage / logging / post-analysis}—so the kernel is not a static rule set but a versioned, deployable, evidence-ready runtime artifact referable from governance. controls-core is not a substitute for the whole control plane, but it supplies a concrete mapping to this chapter’s kernel interfaces, dual-plane split, and Chapter~5’s bounded deployment contract.

\begin{deliverabletable}
Policy engine & schemas + rules + decision traces & deterministically replayable over agreed event types & Runtime team & P1 \\
Input/Output \Guardrails{} & classifiers, validators, rewrite templates & measurable bypass-rate reduction vs baseline & Security ML & P1 \\
Retrieval security service & corpus labeling, citation binding, retrieval allowlists & reduced real-world injection/poison harm & Platform + App & P1 \\
Tool controller & capability matrix, consent, approvals, sandboxing & high-impact actions require explicit policy paths & Agent platform & P2 \\
Audit plane & append-only logs, correlation IDs & evidence can reconstruct events within budget & Platform SRE & P2 \\
Evaluation service & suites, dashboards, regression & versioned packs satisfy release gate & QA / Assurance & P1 \\
\end{deliverabletable}

\keytakeaway{A reference architecture must be readable as both “layers” and “flows”: minimal three-plane fields enable joinable evidence closure; synchronous/asynchronous dual planes define hot vs cold responsibilities; and \SafetyKernel{}’s minimal interfaces make \DecoupledSafety{} an engineering-acceptance boundary for Chapters 4--5.}

% =============================
% Chapter 4
% =============================
\chapterstarter{Runtime Enforcement and Policy Orchestration}{
This chapter specifies the \emph{runtime system layer}. It first separates the \textbf{online hot path} from the \textbf{out-of-band cold path}; then covers risk closure and enforcement for \RAG{} and tools. The core is a \textbf{two-stage fail-safe service protocol} (a report-form compressed notation of T1--T3, D1--D4, S1, Tier-R1/R2), followed by \textbf{ORA (Observe--Route--Adapt)} as a controlled out-of-band adaptation loop aligned with the asynchronous plane in Chapter 3 and the evidence requirements in Chapter 5.
}{
\begin{itemize}
  \item Online vs out-of-band: signal admission and discipline for expensive actions.
  \item Risk closure, policy lifecycle, \RAG{} and tools.
  \item Fail-safe protocol + ORA loop aligned with evaluation evidence fields.
\end{itemize}
}

\section{Runtime instantiation of the computable envelope}
\DecoupledSafety{}. Chapter~2 motivates the \textbf{computable envelope}; here we fix its \textbf{runtime instantiation}: finite triggers (T1--T3), finite actions (D1--D4), invariant S1, recovery tiers (Tier-R1/R2), and joinable evidence fields per decision---not ``many detectors,'' but a labeled protocol whose outputs are enforceable within budget and citeable in evidence packs.

The controls-core execution strategy illustrates implementation as a \textbf{latency-bounded multi-tier cascade}: detectors are partitioned by estimated latency into priority tiers; within a tier execution is parallel; after each tier a cumulative score is computed; with fail-fast, crossing a threshold skips remaining tiers. The lightest path begins with \textbf{language detection, exact matches / compiled regex, structural heuristics}, escalates through \textbf{classic ML / ONNX, small advanced heuristics, relatively fast SLM \& LLM}, and invokes \textbf{LLM-as-Judge} only in borderline bands—compressing open semantics into finite streams, finite escalation paths, and finite triggers.

\section{Online hot path vs out-of-band cold path: division of labor}
\ControlPlane{}. Chapter~3 fixes dual-plane correctness and budgets; here we instantiate triggers T1--T3 and actions D1--D4 under those budgets---only the tiered-timeout and S1 accounting specifics below extend that split.

\section{Risk-adaptive control loop}
\ControlPlane{}. Online enforcement and guarantee closure align as:
\begin{enumerate}
  \item \textbf{Measure.} Offline/online measurement produces decision signals (attack success rate, drift, refusal quality, latency).
  \item \textbf{Profile.} Aggregate risk exposure by model and scenario given deployment context.
  \item \textbf{Select.} Choose control depth under risk/cost/latency tolerance (rules vs light detection vs heavy detection; tool constraints; retrieval segmentation).
  \item \textbf{Enforce.} Deterministically apply actions via \SafetyKernel{} interfaces.
  \item \textbf{Evidence and feedback.} Emit audit records; failures drive suite updates, red-team prioritization, and policy revision.
\end{enumerate}

\section{Policy decision lifecycle}
\DecoupledSafety{}. Per request:
\begin{enumerate}
  \item Extract risk signals (topic sensitivity, channel trust, retrieval scope, tool exposure).
  \item Bind the request to applicable policy packs and versions (tenant, geography, workload class).
  \item Decide intervention: allow, rewrite, constrain, sandbox, escalate, block, degrade—with explicit rationale codes.
  \item Record decisions and replayable evidence payload pointers.
  \item Feed events and near-miss injection attempts into offline suites and shadow evaluation.
\end{enumerate}

\section{Retrieval security (\RAG{})}
\ControlPlane{}. Retrieval brings third-party text into the prompt boundary; treat retrieved chunks as untrusted, code-like content. Controls include segmentation (public vs confidential), metadata and lineage requirements, retrieval-time filtering, chunk-bound citation, contradiction handling, bounded privilege for sensitive-slice answer derivation, and batch/online poisoning detection hooks. Retrieval policy must co-evolve with evaluation suites containing poisoning and injected-document scenarios.

\section{Tool invocation control}
\ControlPlane{}. Tools map language into side effects. Enforce tenant/workflow capability matrices, parameter type checks, dynamic approvals for high-impact actions, sandbox execution when feasible, persistent rate limits, and separation of read-only vs mutating tools. Orchestrators must not create parallel paths that bypass recorded decisions.

\section{Two-stage fail-safe service protocol (report-form compressed)}
\DecoupledSafety{}. This section compresses research-style trigger–action–invariant–recovery into an \emph{acceptance-ready narrative}: implementations may retain internal numbering, but external docs need consistent “stage–trigger class–action class–invariant–recovery tier.” \textbf{Stage 1} completes detection and adjudication inside \SafetyKernel{}; \textbf{Stage 2} migrates user-visible service forms within the declared degradation space and writes evidence. Recovery may extend into the cold path but must respect authority and version discipline.

\subsection{Trigger classes T1--T3 (hot-path readable)}
\begin{itemize}
  \item \textbf{T1 (policy/spec conflict)}—mutually exclusive rules, unauthorized tool intent, policy-version/runtime mismatch, or obligations that cannot be simultaneously satisfied for the current \texttt{event}. The hot path must emit a clear \texttt{rationale\_code}; if not locally resolvable, enter Stage 2.
  \item \textbf{T2 (socio-technical/operational anomaly)}—rate/source/tenant context shifts, dependency degradation, approval queues approaching/breaching \textsc{sla}. The hot path may apply low-cost actions (rate limiting, read-only mode, tool-surface contraction); if S1 remains violated, enter Stage 2 D-actions.
  \item \textbf{T3 (validator uncertainty / insufficient evidence)}—low confidence classifiers, missing retrieval lineage, inability to bind key assertions to evidence-plane fields. Default conservatism: choose D1/D3 or D4 rather than “guess and pass” on the hot path.
\end{itemize}

\noindent\textbf{Finite-state reading.} Implementations should ensure T1--T3 cover the hot-path state space without ``undefined'' outcomes: each normalized \texttt{event} reaches a labeled trigger class (possibly multiple signals mapped to one class), then a Stage~1 decision that either resolves inside \SafetyKernel{} or hands off to Stage~2 with explicit \texttt{rationale\_code}. Uncertainty routes to conservative D-actions or Tier-R rather than silent defaults---this keeps the fail-safe layer a finite labeled transition system aligned with evidence fields.

\subsection{Degradation action space D1--D4 (controlled quadrants)}
When triggers hold, only the following \emph{preconfigured} actions are allowed; each carries business-impact tags and required evidence columns (aligned with Chapter 5 E1--E3):
\begin{itemize}
  \item \textbf{D1}—templated refusal or fixed safe response (predictable, measurable utility cost).
  \item \textbf{D2}—restricted stateless sandbox: disable/shrink high-impact tool sets; contract retrieval slices or index generations.
  \item \textbf{D3}—sanitize/redact model outputs before returning (with citeable rationale codes).
  \item \textbf{D4}—hard-stop the request: return an error code and recovery guidance; do not implicitly retry until budgets are exhausted.
\end{itemize}

\subsection{Synchronous service invariant S1}
\textbf{S1 (synchronous respondability):} any synchronous \textsc{api} must still return a \emph{machine-actionable} response body (including degradation payloads or error codes) within the declared latency budget under degraded conditions; it must not block indefinitely and must return tenant-contract-consistent rationale fields. Violating S1 is an implementation defect, not “the model is not smart enough.” \textbf{Quantitative baseline (contract-dependent):} interactive gateways often declare end-to-end budgets such as $T_{\mathrm{budget}} \le 800\,\mathrm{ms}$ \textsc{p95}; enforcement layers must reserve nested slices for normalization, kernel tiers, and evidence append while preserving S1---otherwise acceptance contracts must shrink scope.

\subsubsection*{S1 budget allocation and tiered cascade timeouts}
\textbf{Hard accounting (informal).} Let $T_{\mathrm{budget}}$ denote the synchronous envelope for one request (e.g.\ gateway \textsc{p95}). The budget is a \textbf{feasibility constraint} on components:
\[
  T_{\mathrm{norm}} + \sum_{k \in \mathrm{tiers}} T_k^{\mathrm{alloc}} + T_{\mathrm{dispatch}} + T_{\mathrm{append}} \;\le\; T_{\mathrm{budget}},
\]
where $T_{\mathrm{norm}}$ covers parsing and correlation; each tier $k$ receives an explicit slice $T_k^{\mathrm{alloc}}$ (parallel intra-tier work still consumes the tier slice); $T_{\mathrm{dispatch}}$ covers controlled downstream actions; $T_{\mathrm{append}}$ covers synchronous evidence-append completion required for governed execution. Tiered cascades must enforce \textbf{nested deadlines}: if tier $j$ exceeds its slice, skip or short-circuit remaining discretionary tiers under fail-fast rules, escalate to conservative D-actions, or refuse---never silently wait beyond the remaining envelope. Cascaded timeouts thus prevent one detector layer from inducing passive S1 breach on the whole path.

\subsubsection*{Finite labeled transitions (informal FSM sketch)}
Hot-path outcomes reduce to a finite partition for observability: states \(\{\mathrm{N},\mathrm{T},\mathrm{D},\mathrm{B}\}\) for nominal, trigger-active (T1--T3 class asserted), degraded-in-D (inside D1--D4), blocked/error-terminal. \textbf{Infrastructure interrupts (hierarchical):} faults such as \texttt{decision\_record\_id} mint failure, mandatory evidence-row loss, or evidence-bus partition are \emph{not} modeled as ordinary matrix inputs alongside T1--T3; they act as \textbf{global interrupts} that preempt policy/trigger evaluation and force \textbf{B} (blocked terminal) or a conservative \textbf{D} envelope with explicit \texttt{rationale\_code}, independent of which logical T-state would otherwise apply. The table below covers \textbf{semantic/policy} signals only.

\begin{center}
\footnotesize
\begin{tabular}{@{}lcccc@{}}
\toprule
 & N & T & D & B \\
\midrule
T1--T3 raised & T & T & T & -- \\
D1--D4 selected & D & D & D & D \\
S1 breach & B & B & B & B \\
\bottomrule
\end{tabular}
\end{center}

\subsection{Heterogeneous recovery tiers Tier-R1 / Tier-R2}
\begin{itemize}
  \item \textbf{Tier-R1 (rule-based recovery)}—state clearing, feature-flag rollback, switching to read-only, reverting to lighter detection paths; completes in seconds to minutes; evidence appends automatically. \textbf{Assurance posture during recovery:} while Tier-R1 steps have not converged to a citeable steady state, synchronous traffic must remain \emph{fail-closed on assurance} relative to Track~A claims---typically conservative refusal or degraded capability modes---not silently bypass gates to restore availability without evidence alignment.
  \item \textbf{Tier-R2 (human/normative recovery)}—restore full capability or elevate privilege only after legal/security/data governance approval; must produce \texttt{approval\_ref} and version-monotone policy logs; runtime execution identities must not self-approve.
\end{itemize}
\textbf{Authority boundary:} who may trigger Tier-R2 and which fields may change must be configured separately from execution platform service accounts.

\section{Controlled out-of-band adaptation loop (ORA: Observe--Route--Adapt)}
\ControlPlane{}. ORA is not “online self-learning.” It is a controlled loop on the \textbf{asynchronous OOB adaptation plane} of Chapter 3: it consumes persisted evidence and aggregated metrics only, produces \emph{next-version} releasable configurations, and enters the hot path only through approval and release gating.

\paragraph{What the online path does.} Within \textsc{sla}: evaluate \texttt{event}s under current \texttt{policy\_version}, dispatch actions, append evidence; when T1--T3 hold, enter D1--D4 and Tier-R1 automations. It \textbf{does not} perform full retraining, global threshold sweeps, full index rebuilds, or unapproved policy rewrites.

\paragraph{What the OOB path does.} Consume Observe aggregates: drift, false-positive/false-negative structure, evolving attack surfaces. In Route, generate candidates (thresholds, retrieval segmentation policies, tool allowlist changes, model-route weights). In Adapt, pass offline evaluation and release gating to ship as versioned artifacts. Outputs are always “new version configuration packs,” never hot-path in-memory mutations.

\paragraph{Signals allowed on the online path.} Single-request decidable and budgeted: rule features, cached scores, light retrieval checks, one pre-warmed heavy check, lineage-missing flags bound to the current \texttt{correlation\_id}. Common property: evaluable within S1 and without cross-request global optimization.

\paragraph{Expensive adaptations that must be asynchronous.} Corpus/index refresh at scale, cross-tenant aggregated tuning, full red-team replays, structured policy DSL rewrites, new detector families, and changes to Track A contract semantics. These must run on the ORA cold path with auditable candidate–evaluation–approval–effective chains to avoid online hidden bifurcation.

\paragraph{Concurrency and policy-version monotonicity.} Multiple ORA streams may propose overlapping configuration candidates concurrently. Candidates must reconcile through a version-control discipline (branch, review, merge) prior to ship: merges must preserve \emph{strict monotonicity} of externally effective \texttt{policy\_version} increments and forbid conflicting edits from overwriting production-effective policy without release gating---Git-like workflows are the operational pattern; ad hoc concurrency without merge rules recreates invisible bifurcation under load.

\begin{figure}[htbp]
  \centering
  \IfFileExists{figures/Fig4.png}{\includegraphics[width=0.92\textwidth]{figures/Fig4.png}}{%
    \fbox{\parbox{0.92\textwidth}{\vspace{22mm}\centering\textbf{Figure 2 placeholder: figures/Fig4.png}\par\vspace{22mm}}}%
  }
  \caption{A dual-timescale risk-adaptive security closure loop}
  \label{fig:dual-timescale-loop}
\end{figure}
\noindent This figure depicts a dual-timescale risk-adaptive control system: the hot path performs instance-level runtime enforcement; the cold path organizes benchmarks, risk profiles, release gating, and regression into an evolving policy and evidence loop. Evaluation is a decision-signal compiler; evidence and release gating elevate judgment and execution into releasable, auditable, regressable guarantee objects.

\evaluationnote{Runtime assurance targets}{Hot path: residual risk, utility impact, p95/p99 latency, S1 violation rate. Fail-safe: T*$\rightarrow$D* transition frequency and error-code distribution. ORA cold path: candidate-to-ship cycle time and rollback count. Cross-validate against Chapter 5 minimum evidence pack fields.}

\noindent\textbf{Engineering correspondence.} The same submodule design shows that hot/cold split is not abstract methodology but an implementation constraint. controls-core lists prompt, response, and dataset as first-class inputs and reserves logging, post-analysis, ML training pipelines, dataset scanning, streaming detector design, strict evaluation \& resource testing protocols, and update mechanisms for the cold path outside the hot path. The system naturally spans two timescales: millisecond-to-second online cascades and fail-fast decisions versus offline adaptation and verification around logs, datasets, policy refresh, resource testing, and deployment configuration. ORA is therefore not an optional optimization in this volume but the cold-path governance complement required for a runtime kernel such as controls-core.

\subsection*{Runtime state \(\rightarrow\) external claim contraction}
Guarantee statements in marketing, contracts, and release summaries must \textbf{tighten in lockstep} with runtime signals (Chapter~5 bounded contracts): when triggers, degradation, S1 breaches, recovery, or ORA shipments change what the system can evidence, external claims must contract or cite a narrower acceptance tier---not lag behind production behavior.

\begin{center}
\footnotesize
\begin{tabular}{@{}p{0.34\textwidth}p{0.58\textwidth}@{}}
\toprule
\textbf{Runtime signal} & \textbf{Contract / claim consequence (illustrative)} \\
\midrule
T1--T3 active & Narrow capability claims; disclose conflict handling; no ``fully safe'' Track~A wording unless fail-safe resolves inside contract \\
D1--D4 entered & Utility and feature claims contract to degraded modes named in acceptance contract \\
S1 breach pattern & Latency/availability claims contract; investigation clause until sustained compliance \\
Tier-R1 mid-recovery & Suspend expanded Track~A claims until citeable steady state (fail-closed on assurance) \\
Tier-R2 / \texttt{approval\_ref} & Elevated capability claims only after signed approval + version bump in pack \\
ORA candidate not shipped & No new external guarantee until Release Evidence Pack updated under new \texttt{policy\_version} \\
\bottomrule
\end{tabular}
\end{center}

\keytakeaway{Chapter 4 clarifies three things: hot-path T/D/S discipline, cold-path Tier-R2 and ORA adaptation, and joinable evidence end-to-end—aligning enforcement protocols with the three-plane model (Chapter 3) and release governance (Chapter 5).}

\noindent\textbf{Dual-timescale closure.} Evaluation acts as a cold-path compiler relative to the hot path: Chapter~3 fixes the dual-plane correctness boundary; Chapter~5 specifies evidence packs; Figure~\ref{fig:dual-timescale-loop} summarizes how versions, suites, and release artifacts connect across timescales.

% =============================
% Chapter 5
% =============================
\chapterstarter{Evidence Engine, Claim Discipline, and Release Defense}{
This chapter upgrades “evaluation” into an enterprise \emph{Evidence Engine} and unifies it with \textbf{claim discipline} and the \textbf{bounded deployment contract} (previously a separate chapter in early drafts). It organizes measurement obligations and evidence retention with an RQ-first structure; uses benchmark constitutions and acceptance contracts to constrain deployable conclusions; expresses incommensurable risk semantics via metric hierarchies and multi-track discipline; and compiles outputs into a unified governance decision board through two evidence packs (a Mechanism Evidence Pack and a Release Evidence Pack) for \emph{release defense}—auditable, replayable, differentially rerunnable, and externally stateable.
}{
\begin{itemize}
  \item RQ-first: question $\rightarrow$ evidence $\rightarrow$ release conclusion (lock the question before the score).
  \item Preserve: benchmark constitution / acceptance contract / metrics hierarchy / multi-track evaluation.
  \item Add: Track A/B claim discipline; Minimum Mechanism Evidence Pack; Minimum Release Evidence Pack.
\end{itemize}
}

\section{Semantic bifurcation with governance convergence}
\EvalFramework{}. In the dual-timescale control system, \EvalFramework{} does not output an abstract total score; it organizes multiple incommensurable semantic spaces in enterprise security. There is no single evaluation semantics: global robustness, jurisdictional compliance, language-specific risk, and organizational obligations differ in test objects, aggregation logic, and release semantics. Multi-track evaluation is therefore not “run more tests,” but \emph{semantic bifurcation}: retain independent semantics per track on a shared execution backbone and achieve \emph{governance convergence} via a unified decision board, acceptance contract, and evidence pack system. The point is not score compression, but unified governance: unify the decision “ship / ship reduced / canary / rollback” and its evidence references, not hiddenly compress distinct semantics into one number.

What matters structurally is unified governance decisions—\emph{whether} to ship, under \emph{which} capability pack, under \emph{which} bounded contract—not forcing tracks into one score. Any cross-track compression must disclose weights, losses, and applicability and be recorded as a new claim premise; it cannot be treated as an inherent “single-number truth.” Evaluation in this chapter is therefore not a scoring tool first but the \emph{semantic compilation layer} of the guarantee production system: it determines in which semantic spaces and with what evidence strength the enterprise may claim what guarantees.

\section{From measurement to control choice: evaluation outputs as decision signals}
\EvalFramework{}. Evaluation outputs do not terminate in reports alone; they are compiled into \emph{decision signals} consumable by \ControlPlane{} and release governance. They answer not “which model scores higher,” but which scenarios require deeper control, which capability packs must contract, which jurisdictions/languages cannot enter the deployable track yet, and which conditions must trigger D1--D4, Tier-R2, or release gating. Evaluation’s core function is to translate multi-track measurement into inputs for \texttt{policy} selection, capability boundaries, and governance actions.

Such a decision signal has governance meaning only when (i) it cites a concrete RQ, track identifier, and acceptance-contract clause; (ii) it can be reconstructed via \texttt{correlation\_id} and \texttt{decision\_record\_id} into concrete events, decisions, and samples; (iii) it can be cited by the release evidence pack—not left on a local dashboard or in an experimental conclusion alone. Only then does measurement graduate from “scores” to governance signals usable in release defense; RQ-first, benchmark constitution, acceptance contract, multi-track discipline, and the two evidence packs in the rest of this chapter are organized around this compilation relation.

controls-core illustrates that decision signals need not derive from a single composite score; they may come from structured intermediates inside a runtime cascade—e.g.\ language-detection filtering outcomes, tier hit/skip counts, cumulative score evolution, fail-fast trigger points, and conditional bands for advanced pipelines and LLM-as-Judge. When these signals are written stably into the evidence graph and bound back to \texttt{policy\_version}, deployment regime, and input object type (prompt / response / dataset), they are not mere debug data but compilable decision signals for control selection and release governance.

\section{Evidence compilation as a governance function}
\EvalFramework{}. \EvalFramework{} organizes benchmark constitution, track discipline, metrics hierarchy, and multi-track suites into a measurement system for risk semantics; it compiles measurements—via decision signals, Mechanism Evidence Pack, and Release Evidence Pack—into governance objects consumable by \ControlPlane{} and release governance. \textbf{Governance accountability} must bind: which obligations enter suites, which conclusions may ship under Track~A, and when failures force bounded-contract contraction. Naming of the owning organization follows enterprise policy; the required capability is an evidence-compilation authority with signing power over acceptance contracts and release packs—not an enumeration of adjacent teams.

\section{Obligation-to-trigger compilation (illustrative chain)}
\EvalFramework{} / \ControlPlane{}. Natural-language obligations become deployable only when compiled into triggers, rationale codes, and citeable evidence fields---closing the semantic gap criticized as ``gray rhetoric.'' Illustrative pipeline:
\begin{enumerate}
  \item \textbf{Obligation (natural language):} ``prevent unauthorized data exfiltration.''
  \item \textbf{Obligation class:} confidentiality / data handling.
  \item \textbf{Policy fields (examples):} \texttt{allowed\_retrieval\_domains}, sensitivity labels, tool destination allowlists.
  \item \textbf{Trigger / event class:} T1 (policy conflict) when a tool attempts export to an unlisted domain or crosses sensitivity boundaries.
  \item \textbf{Runtime decision:} D4 (hard-stop) or D3 (redact), per contract.
  \item \textbf{Rationale code (example pattern):} \texttt{ERR\_PRIV\_EXFIL\_BLOCKED}.
  \item \textbf{Rationale-code dictionary (tenant-facing):} each emitted \texttt{rationale\_code} belongs to a machine-readable dictionary mapping stable codes to jurisdiction- and tenant-appropriate disclosure text (including regulated explanations where required); release and audit consumers must cite codes, not ad hoc prose.
  \item \textbf{Evidence-plane record:} joined row linking \texttt{policy\_version}, tool invocation id, decision record.
  \item \textbf{Claim template (bounded):} ``Under version $V$, high-sensitivity egress via unapproved destinations was blocked; replay fidelity $N/N$ on sampled traces.''
\end{enumerate}
\noindent\textbf{Compilable vs.\ non-compilable obligations.} Deterministic constraints (regexes, allowlists, structured metadata tags) compile cleanly. Open-ended virtues such as ``harmlessness'' or ``maximum helpfulness'' cannot be enforced as unconditional guarantees---they must appear only as \textbf{statistical claims} with explicit uncertainty and scope in claim discipline (below).

\section{RQ-first: organizing evidence from research questions to release conclusions}
\EvalFramework{}. The evidence engine starts from auditable questions, not a score table. Each RQ binds (i) allowed conclusion types, (ii) required evidence sources, and (iii) minimum entry thresholds for the release evidence pack.
\begin{itemize}
  \item \textbf{RQ1 (deployable joint objective):} under declared \textsc{sla} and cost, can the security–utility–latency joint target be satisfied stably? Evidence: metric hierarchy + online shadow comparisons and slices.
  \item \textbf{RQ2 (fail-safe necessity):} under adversarial and socio-technical pressure, are T1--T3 correctly detected and do they converge under D1--D4 and Tier-R discipline without violating S1? Evidence: Chapter 4 trigger/action logs and replays.
  \item \textbf{RQ3 (ORA OOB loop):} is Observe--Route--Adapt’s marginal value reproducible relative to pure semantic blocking, and does expensive adaptation remain out-of-band and version-monotone? Evidence: controlled experiments + drift monitoring + candidate–evaluation–approval–effective chains.
  \item \textbf{RQ4 (real attack surface and release):} under representative contexts and jurisdiction tracks, is residual risk below the acceptance contract? Evidence: multi-track aggregation + signed approvals + replayable evidence entries.
\end{itemize}

\section{Benchmark Constitution: measurement obligations as coverage discipline}
\EvalFramework{}. The constitution states coverage expectations: failure categories, slicing and aggregation, worst-case reporting, and confidence intervals. Its purpose is not one score but the enterprise’s \emph{measurement obligations}. Skipping a failure category must enter the bounded deployment contract (below) as a non-goal or exception with explicit approval.

\section{Acceptance Contract: signable deployment admission conditions}
\EvalFramework{}. The acceptance contract binds scenarios to thresholds and exceptions, turning measurement into deploy-time admission conditions. It is jointly versioned with \texttt{policy\_version}, detector packs, and (if applicable) index generations/model routing, and requires evidence joinability via \texttt{correlation\_id} / \texttt{decision\_record\_id}. It is the semantic anchor of Track A (deployable conclusions): without contract alignment, high scores cannot be extrapolated as deployable claims.

\subsection{Threshold selection logic (CRM sketch)}
Thresholds must be \emph{defensible}, not merely ``chosen by consensus.'' This technical report adopts a constrained risk-minimization framing (informally, CRM): under a stated risk cap $\delta$ fixed by the acceptance contract (e.g.\ upper bound on harm probability or severity class), thresholds should be justified as feasible points on the \textbf{security--utility--latency Pareto surface}---including where they are explicitly derived as minimizers of utility loss $L(U)$ subject to $P(\mathrm{Harm}\mid \theta)\le \delta$. Practice maps into two actionable families:
\begin{itemize}
  \item \textbf{Evidence-constrained thresholds:} anchored to empirical quantities---historical tail failure rates, replay miss rates, join-field completeness---so the threshold inherits operational meaning from telemetry rather than narrative.
  \item \textbf{Optimization-derived thresholds:} from explicit joint targets (security, utility, latency under S1) documented with weights, losses, and rollback rules; changes are treated as \textbf{new contract premises}, not silent retuning.
\end{itemize}
\noindent\textbf{Threshold contractibility (informal).} A threshold is externally contractible only when its selection rule, data sources, and constraints are referenced in the acceptance contract and backed by Track~A evidence---otherwise it remains an internal tuning knob.

\section{Metrics Hierarchy: turning scores into defensible evidence}
\EvalFramework{}. Keep four layers, with governance emphasizing auditability and sliceability:
\begin{itemize}
  \item \textbf{Security.} attack success probability, bypass rates, injection impact, unsafe tool invocation rate, exfiltration attempts (must locate event types and \texttt{rationale\_code}s).
  \item \textbf{Utility.} task success, correctness proxies, calibrated refusal quality and over-refusal rates (must label scenarios/workflows).
  \item \textbf{System.} latency overhead, p95/p99, throughput impact, cost (must separate hot-path vs cold-path costs; hot path constrained by S1).
  \item \textbf{Governance.} trace completeness, replay fidelity, immutability checks, policy-binding correctness, drill-down capability (must join across three planes).
\end{itemize}

\section{Multi-track evaluation into unified governance: convergence without score collapse}
\noindent The point is not “run two suites,” but preserve semantic bifurcation on a shared execution backbone and converge in a unified governance decision board as required by the bounded contract.
\begin{figure}[htbp]
  \centering
  \IfFileExists{figures/Fig2.png}{\includegraphics[width=0.92\textwidth]{figures/Fig2.png}}{%
    \fbox{\parbox{0.92\textwidth}{\vspace{22mm}\centering\textbf{Figure 3 placeholder: figures/Fig2.png}\par\vspace{22mm}}}%
  }
  \caption{Semantic bifurcation with governance convergence in multi-track evaluation}
  \label{fig:semantic-bifurcation}
\end{figure}
\noindent Tracks share one execution backbone; governance convergence produces release decisions and audit artifacts without collapsing distinct track semantics into one aggregate score.
\EvalFramework{}. Incommensurable obligations need not collapse into one score. The deployable pattern is: \emph{shared execution backbone} (run definitions, artifacts, storage, \textsc{ci}/\textsc{cd} hooks) plus \emph{multiple evaluation tracks} differing in suites, aggregation logic, and report templates, converging at a unified board. \textbf{Clarification:} unified governance unifies evidence-referenced decisions, not scores; any compression must disclose weights and losses and be recorded as a new claim premise under claim discipline (below).

\subsection{Track A / Track B claim discipline (no mixing)}
\textbf{Institutional rule:} \textbf{Track A} carries \emph{deployable conclusions} for production release, external compliance statements, and customer contracts; its aggregation and thresholds must align with online observability and must include a release evidence pack. \textbf{Track B} carries \emph{mechanism/component upper bounds and stress tests} for R\&D ranking and architecture improvement; it \emph{must not} be mixed into Track A narratives or compressed into one score. \textbf{Prohibited behaviors:} writing lab upper-bound results into production release summaries without migration arguments; implicitly merging track weights. Violations are evidence-chain breaks; the release gate should reject.

\begin{evaluationcriteria}
\begin{evaluationmatrixtable}
Input/Output \Guardrails{} & bypass rate, harmful miss rate & task completion, user friction & added latency, tail latency & red-team prompts, adaptive attacks \\
Retrieval controls & poisoning hit rate, unauthorized slice access & answer correctness/usefulness & retrieval latency, index cost & poisoned corpora, injected docs \\
Tool invocation controls & unsafe action rate, privilege abuse & workflow success rate & approval latency & agent attacks, tool chains \\
Multi-track / regional suites & track-specific thresholds & comparable utility baselines & runtime cost, ops effort & localized suites + shared backbone \\
Audit / evidence plane & trace completeness, replay success & time-to-root-cause & storage/log volume & replay drills, sampling audits \\
\end{evaluationmatrixtable}
\end{evaluationcriteria}

\section{Evaluation protocol: offline, online, regression, and release (as an evidence pipeline)}
\EvalFramework{}. \textbf{Offline} suites drive iteration and block defects; \textbf{shadow} deployments compare policy variants under controlled risk; \textbf{red teaming} probes adaptive adversaries; \textbf{regression} gates enforce monotonicity constraints on standard packs when policy/model/detector changes. Release standards combine statistical thresholds, worst-case bounds, required artifacts, and operational approvals, and require artifacts to be joinable back to replayable event samples via \texttt{policy\_version} and evidence keys.

\section{Minimum Mechanism Evidence Pack}
\EvalFramework{}. The mechanism pack answers “is the mechanism actually working”: interface closure, replay reproducibility, hot-path budget control. It primarily supports internal reviews and Track B, but also underpins Track A’s release pack.

For a layered runtime kernel, the Mechanism Evidence Pack can be further itemized as evidence that hot-path execution discipline actually occurred; the controls-core implementation anchor is in the Controls-core subsection of Chapter~3; this chapter does not restate its cascade structure.
\begin{itemize}
  \item \textbf{Interface closure and field completeness.} Missing-rate of minimal three-plane fields in production and replay datasets; join success rate of \texttt{event}$\rightarrow$\texttt{decision}$\rightarrow$\texttt{evidence} via \texttt{correlation\_id}/\texttt{decision\_record\_id}.
  \item \textbf{Replayability.} Under the same \texttt{correlation\_id} in differential reruns: consistent \texttt{policy\_version}, stable \texttt{rationale\_code}, and reproducible key decisions (allow bounded variation only in explicitly labeled nondeterministic intervals).
  \item \textbf{Fail-safe unit coverage.} Replay coverage for T1--T3 triggers and D1--D4 migrations; S1 violation detection and alert chain; Tier-R2 authority-boundary checks (no self-approval by runtime identities).
  \item \textbf{Budget and cost discipline.} Hot-path \texttt{deadline\_budget\_ms} hit rate, timeout distribution, degradation hit rate; cold-path ORA queue delays and rollback frequency.
\end{itemize}

\section{Minimum Release Evidence Pack (Release Evidence Pack)}
\EvalFramework{}. The release pack supports release defense: unify decisions while preserving track semantics and bounded contracts. Suggested minimal structure:
\begin{itemize}
  \item \textbf{(A) Version and scope metadata:} \texttt{release\_id}, \texttt{policy\_version}, versions of model routing/detectors/index generations; applicable tenants/jurisdictions/languages/workflows; allowed tool surface and retrieval domain.
  \item \textbf{(B) Contract excerpt:} acceptance-contract thresholds, exceptions, S1, and allowed D1--D4 / Tier-R2 conditions.
  \item \textbf{(C) Track A evidence:} sliced reports over security/utility/system/governance metrics (worst-case and confidence intervals), aligned with online shadow or historical baselines.
  \item \textbf{(D) Fail-safe evidence (E1--E3):} T1--T3 coverage and correctness replays; D1--D4 migration frequency, error-code distribution, p95/p99; S1 violation rate; Tier-R2 approval samples (\texttt{approval\_ref}) and authority boundary checks.
  \item[] \textbf{Deployment variants and configuration-governance bridge (if applicable).} If the runtime kernel uses capability-pack tiers and deployment variants, the Release Evidence Pack should record that the \emph{actually enabled} form matches the boundary referenced by the contract. The controls-core example (capability packs and deployment variants) is in the Controls-core subsection of Chapter~3; this chapter does not expand engineering detail inside the evidence-pack template.
  \item \textbf{(E) ORA evidence (E4--E6, if applicable):} drift monitoring and false-positive costs; auditable candidate–evaluation–approval–effective chain samples; expensive-adaptation async discipline and rollback records.
  \item \textbf{(F) Governance artifacts:} risk acceptance/exceptions, rollback plan, monitoring thresholds, audit sampling plan.
  \item \textbf{(G) Unified decision output:} ship / ship reduced / canary / rollback, each referencing the evidence items above (not replaced by a single score).
\end{itemize}
\noindent The release pack and claim discipline (below) form the minimal externally stateable set: missing items weaken claim strength or require capability contraction.

\section{Typed release packs, admissibility rules, and anti-abuse tests}
\EvalFramework{}. Treat each Release Evidence Pack as a \textbf{typed artifact}: every section must carry machine-checkable metadata (track, policy version, scope, approval pointers). Subjective narrative alone is inadmissible for Track~A.

\subsection{Admissibility rules (release gating)}
\begin{itemize}
  \item \textbf{Rule 1 (provenance):} evidence whose source is exclusively Track~B, ad-hoc scratch tests, or unaudited one-off runs is \textbf{inadmissible} for Track~A deployable conclusions unless accompanied by an explicit migration argument and re-execution under the shared backbone.
  \item \textbf{Rule 2 (closure):} an entry is admissible only if minimal three-plane join fields are present and replay succeeds within declared budgets; partial screenshots or aggregate dashboards without join keys fail closure.
  \item \textbf{Rule 3 (version alignment):} each cited artifact’s \texttt{policy\_version} must match the release object under consideration; mismatches invalidate the pack for external statements.
\end{itemize}

\subsection{Release gate invariant (organizational intent)}
The governance target is that no Track~A ship decision relies on packs containing unversioned blobs, mismatched \texttt{correlation\_id} samples, or Track~B-only proofs. Operationally, implement an automated \textbf{release-pack validator} in \textsc{ci}/release pipelines (hard failure on structural violations) rather than relying solely on manual vote---manual steps remain for risk acceptance, not for bypassing join closure.

\subsection{Cryptographic and systemic immutability (recommended baseline)}
\EvalFramework{}. To resist retrospective tampering by compromised execution identities, \texttt{decision\_record\_id} issuance and evidence-plane appends should rely on an append-only or hash-chained store (or equivalent tamper-evident log service): each record binds predecessor hash, \texttt{policy\_version}, and partition key so reordering or deletion is detectable. Release gates reject packs whose cited rows fail integrity verification against the declared anchor. \textbf{Discipline:} tamper-evident storage is a \emph{recommended baseline} for strong audit defense; whether it is \emph{mandatory} for a given Track~A claim class depends on deployment regime and acceptance contract---black-box or degraded-observability tiers may cite weaker integrity artifacts only if explicitly bounded there.

\subsection{Ring-fencing external drift without policy bumps}
\EvalFramework{} / \ControlPlane{}. Stateful tools and externally hosted retrieval indexes can shift behavior without a local \texttt{policy\_version} bump. Contracts must either (i)~pin endorsed tool endpoints, corpus generations, and embedding generations explicitly in the acceptance contract and evidence pack, or (ii)~treat unmanaged drift as Tier-B upper-bound evidence until re-validated under Track~A---preventing silent guarantee creep.

\subsection{Anti-abuse: marketing over-claim test suite}
Maintain a regression suite that ingests claims from release notes or external blurbs and attempts to \textbf{compile each sentence} to Track~A evidence pointers. If compilation lands only on Track~B mechanisms or lacks joinable rows, the gate fails with an ``over-claim'' classification---preventing aspirational language from shipping without evidence lineage.

\section{Claim discipline}
\noindent This section specifies \textbf{under this framework} which safety statements constitute valid enterprise guarantee statements. Any claim in external materials, contract text, release notes, or executive summaries must be constrained to explicit scope, premises, degradation semantics, deployment regime, and evidence-pack references. Any statement that cannot be compiled into this form must not be presented as a deployable guarantee; it must be downgraded to research agenda, design intent, or an explicit non-goal.

\subsection{Claim type system (type-checking against the evidence regime)}
Beyond prose templates, assign each candidate statement a \textbf{claim type} keyed to the deployment regime (white-box / semi-white-box / black-box): mechanism-replay claims, statistical tail-risk claims, S1-only claims, capability-boundary claims, etc. A statement is \emph{well-typed} only if the evidence regime listed in its guarantee grammar actually supports the evidentiary constructors it uses---for example, black-box regimes cannot legitimately carry mechanism-level replay assertions.

\subsection{Guarantee grammar (required fields): what statements are allowed}
A guarantee statement that is allowed into external materials or release defense must include at least:
\begin{enumerate}
  \item \textbf{Scope:} applicable tenants, jurisdictions, languages, workflows, and capability-pack boundaries;
  \item \textbf{Invariants:} properties the system commits to maintain (e.g., S1, allowed action space, authority boundaries);
  \item \textbf{Degradation semantics:} the minimum service form entered when premises break (e.g., D1--D4, Tier-R1/R2, capability contraction);
  \item \textbf{Preconditions:} observability, governance, dependencies, and contract premises required for the guarantee to hold;
  \item \textbf{Evidence basis:} citeable evidence-pack items, track IDs, replay keys, versions, and approvals;
  \item \textbf{Deployment regime:} whether the guarantee holds under white-box, semi-white-box, or black-box conditions.
\end{enumerate}
\noindent Missing any item invalidates the claim for release defense and external statements.

\subsection{Claim-to-evidence compilation}
Each deployable claim must compile into acceptance contract clauses, release evidence pack references, and a release decision object (ship / ship reduced / canary / rollback).

\subsection{Non-goals}
This technical report does not claim:
\begin{enumerate}
  \item global, complete, decidable safety proofs for arbitrary tasks, tools, and contexts;
  \item that a single offline red-team run, a single attack success rate, or a local benchmark score equals production safety;
  \item that Track B or any mechanism/upper-bound track results can be used as deployable conclusions without explicit migration arguments;
  \item that ``a smarter model'' substitutes for policy, evidence, contracts, and release objects;
  \item that a system is ``audited/replayable/defensible'' without citeable \texttt{policy\_version}, \texttt{correlation\_id}, \texttt{decision\_record\_id}, or approval records.
\end{enumerate}

\subsection{Bounded guarantees}
Under explicit premises, allowed bounded guarantees include synchronous interfaces satisfying \textbf{S1} under degradation; when \textbf{T1--T3} hold, degradation only within \textbf{D1--D4}; Tier-R2 boundaries institutionalized; joinable three-plane fields; and replayable ORA when enabled.

\subsection{Preconditions}
Bounded guarantees depend on semantic premises (decidable taxonomy/schema), deployment/observability premises (immutable logs, joinable fields), governance premises (Tier-R2 chains), and dependency premises (tools/retrieval availability calling T2). If premises fail, claims contract.

\noindent\textbf{Engineering preconditions (supplement).} Cover configuration, logging, storage, and update chains (controls-core anchor in Chapter~3).

\paragraph{\texttt{evidence\_schema\_version} backward compatibility and migration admissibility.} Regression and differential reruns must declare a compatible \texttt{evidence\_schema\_version}: within a major, new minors are additive and readers for that major must accept prior minors; a new major requires explicit migration clauses in the acceptance contract and coordinated updates to validators and Release Evidence Pack templates---otherwise historical packs cannot join to current audits. \textbf{Migration admissibility:} a schema major bump is \textbf{inadmissible} for uninterrupted Track~A continuity until reruns pass under the migration plan and release-pack validators are updated---otherwise release must contract scope or treat the transition as a signed exception with replay proof.

\subsection{Degradation regimes}
When premises fail: contract capability packs first; conservative D1--D4 next; Tier-R1/R2 with evidence; synchronize external claim contraction.

\subsection{Deployment regime layering (white-box / semi-white-box / black-box)}
\begin{itemize}
  \item \textbf{White-box:} full three-plane fields; differential reruns; stronger release defense.
  \item \textbf{Semi-white-box:} partial fields with disclosed approximation boundaries.
  \item \textbf{Black-box:} only weaker guarantees (capability boundaries, S1, error codes).
\end{itemize}

\subsection{Minimal claim forms (external statements)}
\begin{enumerate}
  \item \textbf{Scope sentence:} Within \{tenant/jurisdiction/language/workflow/capability pack\} scope, under \{\texttt{policy\_version}\}, the system provides \{bounded guarantee\}.
  \item \textbf{Invariant sentence:} For synchronous interfaces, the system satisfies \{S1\}; when \{T1--T3\} hold, it degrades only within \{D1--D4\}.
  \item \textbf{Evidence sentence:} The guarantee is supported by \{Release Evidence Pack items / \texttt{correlation\_id} / \texttt{decision\_record\_id} / \texttt{approval\_ref}\}, replayable and citeable.
  \item \textbf{Degradation sentence:} When premise \{list\} fails, the system enters \{minimum service form\} and contracts external claims accordingly.
\end{enumerate}

\section{Bounded deployment contract}
\textbf{Synthesis.} A \textbf{bounded deployment contract} is the \textbf{signed instantiation} of the claim forms above for a concrete deployment: it binds what may be stated externally to explicit scope, premises, and evidence-pack references.

\section{Evidence graph and release governance}
\EvalFramework{} / \ControlPlane{}. Release decisions cite an evidence graph that links versions, suites, slices, and approvals into a joinable release narrative. Unified governance decision \(\neq\) unified score.

\paragraph{Operating cadence.} Minimum enterprise cadence and roadmap appear in Appendix~\ref{app:rd-framework}.

\researchagenda{Evaluation frontier}{How can organizations jointly optimize security, utility, cost, latency, and auditability under measurable confidence intervals—while keeping multi-track evidence stable as bounded contracts in release defense and long-term audit?}

\keytakeaway{Chapter~5 unifies the Evidence Engine with claim discipline and bounded deployment: multi-track semantics without score collapse; Track A/B separation; Mechanism and Release Evidence Packs; typed admissibility with tamper-evident joins where applicable; signable contracts and governance graphs---closing guarantee grammar with release defense in one chapter. Operating rhythm sits in Appendix~\ref{app:rd-framework}.}

% =============================
% Chapter 6 (formerly Chapter 7)
% =============================
\chapterstarter{Jurisdictional Semantics, Language Gates, and Boundary Intelligence}{
\textbf{Primary scientific object:} \emph{jurisdictional semantic compilability}---which obligations and languages admit Track~A deployable conclusions under explicit suites and contracts. \textbf{Chinese and multilingual AI safety} serve as the canonical rigorous case study for dense linguistic and regulatory boundaries (tokens, refusal/over-refusal semantics, localized suites); the same discipline applies to EU AI Act, HIPAA, or sector regimes when encoded as taxonomy $\rightarrow$ constitution $\rightarrow$ evidence rows.
}{
\begin{itemize}
  \item Encode jurisdictional semantics into taxonomy and acceptance contracts.
  \item Language gates decide Track~A vs Track~B eligibility per jurisdiction.
  \item Regional detectors and multilingual suites produce release-evidence artifacts with shared join keys.
\end{itemize}
}

\section{Why geo-specific semantics is core intelligence}
\ControlPlane{} / \EvalFramework{}. Language and jurisdiction differences are not “presentation-layer differences,” but differences in obligation ontology, risk semantics, explainability requirements, data sovereignty boundaries, and acceptable residual risks. Treating them as auxiliary creates fractures: \textbf{policy fracture} (global-average obligations misalign with local deployable conclusions), \textbf{evidence fracture} (decision records lack local audit fields), and \textbf{release fracture} (Track A deployable conclusions end up citing Track B upper bounds or translated metrics, yielding indefensible narratives). Therefore multilingual and China intelligence must enter the main evidence engine and bounded contract.

\section{Semantics and obligation baselines: from regulatory classes to testable obligations}
\EvalFramework{}. Compile geo-specific obligations into benchmark constitutions and acceptance contracts:
\begin{itemize}
  \item \textbf{Taxonomy and mapping.} build taxonomy from local regulatory/industry obligations and map into \texttt{policy} constraint fields, thresholds, and exception clauses (into acceptance contracts).
  \item \textbf{Rationale-code dictionary.} harden local semantic differences into \texttt{rationale\_code} dictionaries and explanation templates so refusals/degradations remain replayable and auditable across languages.
  \item \textbf{Fail-safe alignment.} when local obligations conflict with available evidence, trigger Chapter 4 T1--T3 and degrade under D1--D4 and Tier-R disciplines; “language unsupported” must not be a bypass without evidence.
\end{itemize}

\noindent\textbf{Compliance-to-evaluation compilation.} This is effectively a compilation pipeline: obligation semantics $\rightarrow$ evaluation obligations $\rightarrow$ releasable evidence, producing suite coverage, aggregation logic, rationale dictionaries, detector artifacts, and finally Track A contract and evidence-pack items.

\section{Jurisdictional admissibility for Track A}
\EvalFramework{}. Regulatory and locale-specific obligations define which sentences may appear in Track~A conclusions for that geography. \textbf{Jurisdictional admissibility} means: the obligation taxonomy and acceptance-contract clauses for a jurisdiction are satisfied by suite coverage, rationale-code dictionaries, and detector artifacts \emph{before} the language enters the deployable track. If a jurisdiction’s obligations are only partially compiled, Track~A must contract scope (capabilities, workflows, languages) rather than extrapolate global scores---mirroring admissibility logic in Chapter~5 but at the semantics layer of geo-specific guarantees.

\section{Language gates: entry discipline for Track A}
\EvalFramework{}. Language gates convert deployable-conclusion eligibility into auditable admission. \textbf{Inputs:} suite coverage, false-positive cost, joinable three-plane completeness, hot-path S1 budget impact.

\paragraph{Gate output schema (institutional).} Each gate emits \texttt{language\_gate\_status} \(\in \{\texttt{TrackA\_admissible},\,\texttt{TrackB\_only},\,\texttt{capability\_contracted}\}\) plus, when not \texttt{TrackA\_admissible}, explicit \textbf{capability contraction} (e.g., tools disabled, retrieval domains narrowed, workflows excluded) recorded for the bounded deployment contract. \texttt{TrackB\_only} forbids Track~A marketing or compliance sentences for that language/jurisdiction until remediating evidence ships; \texttt{capability\_contracted} allows partial Track~A only for the contracted surface.

\section{Multilingual evaluation as first-class tracks}
\EvalFramework{}. Multilingual/regional tracks are first-class: suite versions, aggregation logic, and report templates per track; semantics converge at the governance board under Chapter~5 semantic bifurcation (\textbf{unified decision \(\neq\) unified score})---this section adds only jurisdictional slicing.

\section{Chinese detector development: from semantics to deployable artifacts}
\ControlPlane{}. Core outputs are versioned artifacts:
\begin{itemize}
  \item \textbf{Detector/rule packs.} develop detectors for Chinese/multilingual attack surfaces (tokenization traps, mixed-language injection, regional corpus poisoning) and jointly version them with \texttt{policy\_version}.
  \item \textbf{Evidence binding.} detector outputs must be written into evidence-plane fields (joinable by \texttt{decision\_record\_id}) and referenced in the Chapter 5 release evidence pack.
  \item \textbf{Cold-path iteration.} expensive adaptations (global threshold sweeps, detector training, suite expansion) must happen in ORA cold path with auditable candidate–evaluation–approval–effective chains to prevent online hidden bifurcation.
\end{itemize}

\keytakeaway{Chinese and multilingual capability exemplifies jurisdictional and linguistic boundary intelligence for the control plane and evidence engine: geo-specific obligation mappings, language gates, detector artifacts, and releasable evidence chains so Track~A deployable conclusions remain auditable, replayable, and defensible. The same framework extends to other high-specificity jurisdictional regimes (e.g., sectoral medical rules) when encoded with the same discipline.}

% =============================
% Appendix
% =============================
\appendix

\chapter{Unified R\&D Framework and Operating Program}\label{app:rd-framework}
This appendix specifies operating cadence, gates, and conflict handling for the unified R\&D program.

\section{Theory--system--benchmark--release: turning claims into a running pipeline}
The unified framework is a single traceable chain with feedback---not four silos. \textbf{Theory} fixes claim boundaries; \textbf{System} instantiates \texttt{event}/\texttt{decision}/\texttt{evidence} and runtime protocols; \textbf{Benchmark} produces constitutions, contracts, suites, packs; \textbf{Release} outputs ship decisions. \textbf{Constraint:} externally stateable conclusions cite aligned theory, objects, benchmark evidence, and contract clauses.



\section{Unified security program (single operating grammar)}
The enterprise wraps theory, system, benchmark, and release into one operating grammar: units do not define incompatible object languages, evidence gates, or release semantics. Shared objects include \texttt{policy\_version}, \texttt{correlation\_id}, \texttt{decision\_record\_id}, track discipline, acceptance contracts, and release evidence packs.

Therefore outward-facing outcomes state \emph{under explicit premises, which bounded guarantees hold and how packs support them} rather than ``teams shipped many modules.'' Evidence compilation responsibilities attach to the governance organization named where relevant in Chapter~5.

\paragraph{Internal factorization.} F1 semantics and obligation generation; F2 runtime control implementation; F3 measurement and evidence compilation; F4 release and supervision closure. Engineering anchors appear in the Controls-core subsection of Chapter~3.

\section{How to operate coherently: avoid parallel, fractured evolution}
\begin{itemize}
  \item \textbf{First, one source-of-truth object set.} Unified event, evidence, and version semantics (e.g.\ \texttt{policy\_version}, \texttt{correlation\_id}, \texttt{decision\_record\_id}, tracks, contract clauses). Non-joinable objects break closure.
  \item \textbf{Second, one release gate.} Deployable conclusions carry a Release Evidence Pack; Track~B results do not silently leak into release summaries.
  \item \textbf{Third, one evolution cadence.} Incidents, regression gates, and version-monotone releases co-evolve theory, system, benchmark, and release.
\end{itemize}

\noindent\textbf{Engineering coupling.} Detector expansion, configuration toggles, deployment trimming, dataset scanning, resource testing, and update mechanisms span runtime and evidence layers; disjoint cadences lose a single \texttt{policy\_version}, evidence chain, and release gate.

\section{Operating cadence (minimum cadence)}
\begin{itemize}
  \item \textbf{Each release:} differential reruns and citation checks against the Chapter~5 Release Evidence Pack checklist.
  \item \textbf{Weekly:} review risk, false positives/negatives, fail-safe migrations (T$\rightarrow$D); archive incidents under \texttt{correlation\_id}.
  \item \textbf{Monthly:} update benchmark constitutions and suites; revise acceptance-contract thresholds and expiry policy.
  \item \textbf{Quarterly:} replay drills; verify three-plane join, immutability checks, Tier-R2 boundaries.
\end{itemize}

\section{Roadmap (capability maturity and research upgrade, compressed)}
\begin{itemize}
  \item \textbf{Stage 1: baseline envelope.} Input/output \Guardrails{} + structured audit + minimal acceptance contract and regression gate.
  \item \textbf{Stage 2: context risk.} Retrieval/tool mediation + stronger evidence correlation + shadow evaluation.
  \item \textbf{Stage 3: unified control plane.} Cross-model policy abstraction + automated differential reruns.
  \item \textbf{Stage 4: controlled adaptation.} ORA loop + version-monotone updates + rollbackable expansion of the guarantee boundary.
\end{itemize}

\noindent\textbf{Conflict resolution when release rejects a launch.} The governance board blocks deployment until (i)~contract premises are tightened or capability packs contracted, (ii)~missing evidence-pack rows are remediated with replay proof, or (iii)~an explicit exception with expiry and monitoring riders is signed---otherwise the bounded guarantee cannot be stated.

\chapter{Terminology and responsibility boundaries}
\textbf{Security judgment providers} are responsible for policy interpretation, risk scoring, detector artifacts, explainability payloads, and structured evidence suitable for audits. \textbf{Execution platforms} provide networking, scaling, orchestration \textsc{sla}, and load management. Clear boundaries avoid attribution vacuums where neither side can produce replayable traces. For external communication, internal project names should be mapped to these neutral roles.

\end{document}

